\begin{helper}
Let \(V\) be a finite nonempty set and let \(A\subseteq V\). Then the centered indicator
\[
  v\mapsto 1_A(v)-\frac{|A|}{|V|}
\]
has coordinate sum \(0\).
\end{helper}
\begin{proof}
The sum of the ordinary indicator \(1_A\) over all vertices is \(|A|\), since exactly the elements of \(A\) contribute \(1\). The constant function with value \(|A|/|V|\) has total sum
\[
  |V|\cdot \frac{|A|}{|V|}=|A|,
\]
because \(V\) is nonempty and hence \(|V|\ne 0\). Therefore
\[
  \sum_{v\in V}\left(1_A(v)-\frac{|A|}{|V|}\right)
  =
  \sum_{v\in V}1_A(v)-\sum_{v\in V}\frac{|A|}{|V|}
  =
  |A|-|A|=0.
\]
\end{proof}
